\renewcommand{\Anno}{2026}

\renewcommand{\NumeroProblema}{2}

\renewcommand{\Testo}{
Un camion deve attraversare un deserto lungo 1000~km e consuma 1~litro di carburante per ogni km percorso. Il camion ha un serbatoio con capacità di $C$~litri ed ha la possibilità di creare dei depositi intermedi di carburante lungo il percorso, versandovi parte del carburante residuo del serbatoio. Non si mettono limiti al numero di viaggi che il camion può effettuare.

\begin{enumerate}
\renewcommand{\labelenumi}{(\alph{enumi})}

\item Consideriamo il caso in cui $C=750$ e aggiungiamo il vincolo che il camion possa creare un solo deposito intermedio. Sotto tali ipotesi, determinare la strategia migliore per compiere l'attraversamento del deserto con il minimo consumo totale di carburante, e la quantità di carburante consumata.

\item Nel caso $C=500$, determinare il minimo numero di depositi intermedi di carburante necessari per poter attraversare il deserto.

\item Determinare, in funzione di $C$, il minimo numero di depositi intermedi di carburante necessari per poter attraversare il deserto.

\end{enumerate}
}

\renewcommand{\Risposta}{}

\renewcommand{\Aiutino}{}

\renewcommand{\SoluzioneUfficiale}{}

\renewcommand{\SoluzioniAlternative}{}

\renewcommand{\Osservazioni}{}

\renewcommand{\Commenti}{}
